\section{Nonclosed equations on cuspidal products}
\label{sec:nonclosed-cuspidal-products}\index{Unibranch space!nonclosed cuspidal equations}

A cuspidal normalization can turn a degenerate pulled-back Hermitian form into a positive reference by a smooth potential. On the product below, the nonclosed perturbation has all powers $\ddc$-closed. This gives a fixed mass and an energy comparison even though the background is not closed. A quantitative local equation supplies a positive solution against which every bounded competitor can be compared.

Fix $n\geq2$. Let $E_j=\C/(\mathbb Z+\mathrm i\mathbb Z)$ and $\kappa_j=\frac{\mathrm i}{2\pi}\mathrm dz_j\wedge\mathrm d\bar z_j$. Replace a neighborhood of a point $p\in E_1$ by the cusp with normalization $z\mapsto(z^2,z^3)$; equivalently, the holomorphic functions at the marked point are the functions on $E_1$ whose first derivative there vanishes. Denote the resulting compact curve by $C$ and its cusp by $q$. Put
\[
 X=C\times E_2\times\cdots\times E_n,\qquad
 Y=E_1\times\cdots\times E_n.
\]
The normalization $\pi:Y\to X$ is finite, bijective and a homeomorphism, with generic degree one. The space $X$ is reduced, irreducible and unibranch. Its singular locus is $\{q\}\times E_2\times\cdots\times E_n$.

Choose a Hermitian form $h_C$ on $C$ which near $q$ is the restriction of $\ddc(|x|^2+|y|^2)$. Its pullback is $\eta=\rho\kappa_1$, where $\rho$ is smooth, positive off $p$, and
\begin{equation}\label{eq:cuspidal-reference-density}
 \rho=4|z|^2+9|z|^4\quad\hbox{near }p.
\end{equation}
Write $\lambda$ for its positive average on $E_1$, and set
\[
 h_0=h_C+\sum_{j=2}^n\kappa_j,\quad
 H_0=\pi^*h_0=\eta+\sum_{j=2}^n\kappa_j,\quad
 K_0=\lambda\kappa_1+\sum_{j=2}^n\kappa_j,\quad \Omega=K_0^n.
\]
Forms from the factors are pulled back without further notation.

\begin{proposition}\label{prop:cuspidal-nonclosed-repair}
There are a smooth real mixed form $\alpha$ on $X$ and a smooth real function $v_0$ on $Y$ such that, for every sufficiently small nonzero real $\varepsilon$, the Hermitian form $h=h_0+\varepsilon\alpha$ is nonclosed and
\[
 \begin{gathered}
 H=\pi^*h,\qquad k=H+\ddc v_0=K_0+\varepsilon\pi^*\alpha\geq\tfrac12K_0,\\
 \ddc(H^r)=\ddc(k^r)=0\quad(r\geq0).
 \end{gathered}
\]
Moreover $\int_YH^n=\int_Yk^n=:M>0$.
\end{proposition}
\begin{proof}
Choose real smooth $A$ on $C$, supported away from $q$, and real smooth $B$ on $E_2$. Let
\[
 \alpha=\partial A\wedge\bar\partial B+\bar\partial A\wedge\partial B
        =\partial(A\bar\partial B)+\bar\partial(A\partial B).
\]
The form is real and $\ddc\alpha=0$. At a point choose $\partial A=0$, $\partial\bar\partial A\ne0$ and $\partial B\ne0$. Then
\[
 \partial\alpha=-\partial A\wedge\partial\bar\partial B
                    +\partial\bar\partial A\wedge\partial B\ne0.
\]
Thus $\mathrm d\alpha\ne0$. It vanishes near the singular locus. Its operator norm relative to $H_0$ is bounded on its support, and its norm relative to $K_0$ is bounded everywhere. Choose $|\varepsilon|$ so that both norms after multiplication are less than $1/2$. This proves ambient positivity of $h$, semipositivity of $H$, and positivity of $K_0+\varepsilon\pi^*\alpha$.

The mean-zero equation $\partial_z\partial_{\bar z}v=\lambda-\rho$ on $E_1$ has a smooth real periodic solution by division of its nonzero Fourier coefficients. Its pullback $v_0$ satisfies $\ddc v_0=(\lambda-\rho)\kappa_1$, proving the formula for $k$. The forms $H_0,K_0$ are closed. Both $\mathrm dk$ and $\dc k$ therefore come from the complex surface $E_1\times E_2$. Consequently $\mathrm dk\wedge\dc k=0$ by degree, and $\ddc k=0$. The same identities hold for $H$. The Leibniz rule now gives $\ddc(k^r)=\ddc(H^r)=0$ for all $r$. This uses both torsion identities, rather than pluriclosedness alone. Expanding $(H+\ddc v_0)^n$ and integrating by parts makes every nonconstant term zero. Hence the masses agree; positivity of $k$ makes them positive.
\end{proof}

Fix an integer $s\geq n+3$ before making a further choice of $\varepsilon$. Use the Fourier Sobolev norm
$\|w\|_{H^t}^2=\sum_{m\in\mathbb Z^{2n}}(1+|m|^2)^t|\widehat w(m)|^2$, and write $H^t_0$ for its real zero-$\Omega$-mean subspace. In the remaining formulas $\alpha$ denotes its pullback to $Y$.

\begin{theorem}\label{thm:cuspidal-quantitative-equation}
For sufficiently small nonzero $\varepsilon$, depending on $n,s$ and the fixed forms, there are explicit constants $R,\delta>0$ such that every $g\in H^s_0$ with $\|g\|_{H^s}\leq\delta$ admits a unique $w\in H^{s+2}_0$ with $\|w\|_{H^{s+2}}\leq R$ satisfying
\[
 (k+\ddc w)^n=(f_\varepsilon+g)\Omega,
 \qquad f_\varepsilon=\frac{k^n}{\Omega},\qquad
 k+\ddc w\geq\tfrac14K_0.
\]
The constants may be chosen as in \eqref{eq:cuspidal-contraction-radius}. In particular $w$ is $C^2$.
\end{theorem}
\begin{proof}
Set $b_1=\lambda$, $b_j=1$ for $j\geq2$. The operator
\[
 L_0w=\frac{nK_0^{n-1}\wedge\ddc w}{\Omega}
     =\sum_j b_j^{-1}\partial_{z_j}\partial_{\bar z_j}
\]
has Fourier eigenvalues $\ell(m)=-\pi^2\sum_jb_j^{-1}(m_{2j-1}^2+m_{2j}^2)$. It has an inverse $H^s_0\to H^{s+2}_0$ of norm
\[
 A_0=\sup_{m\ne0}\frac{1+|m|^2}{|\ell(m)|}
       \leq\frac{2\max_jb_j}{\pi^2}.
\]
Let $L_\varepsilon w=nk^{n-1}\wedge\ddc w/\Omega$. It preserves zero mean because $\ddc(k^{n-1})=0$. Fix $\varepsilon_0\leq1$ within the positivity range. Its finite expansion gives
\[
 L_\varepsilon-L_0=\sum_{j=1}^{n-1}\varepsilon^jP_j,
 \qquad P_jw=n\binom{n-1}{j}
    \frac{K_0^{n-1-j}\wedge\alpha^j\wedge\ddc w}{\Omega}.
\]
Choose multiplication bounds $E_j$ for $\|P_j\|_{H^{s+2}\to H^s}$ and set $B=\sum_j\varepsilon_0^{j-1}E_j$. Impose $|\varepsilon|\leq\varepsilon_0$ and $|\varepsilon|A_0B\leq1/2$. The Neumann series gives $\|L_\varepsilon^{-1}\|\leq A:=2A_0$.

The remainder is
\[
 Q(w)=\sum_{r=2}^n\binom nr\frac{k^{n-r}\wedge(\ddc w)^r}{\Omega}.
\]
For $s>n$, Fourier convolution and Cauchy--Schwarz use $\sum_m(1+|m|^2)^{-s}<\infty$ to make $H^s$ an algebra. Smooth coefficient multiplication and the finite expansion therefore give a constant $C$ such that, on the unit $H^{s+2}$ ball,
\[
 \begin{aligned}
 \|Q(w)\|_{H^s}&\leq C\|w\|_{H^{s+2}}^2,\\
 \|Q(w)-Q(z)\|_{H^s}&\leq C(\|w\|_{H^{s+2}}+\|z\|_{H^{s+2}})\|w-z\|_{H^{s+2}}.
 \end{aligned}
\]
Each term has zero integral: for smooth $w$, its integral is $\int w\,\ddc(k^{n-r})\wedge(\ddc w)^{r-1}=0$, and Sobolev approximation gives the general case. Sobolev embedding also gives $D$ with $\|\ddc w\|_{L^\infty,K_0,\mathrm{op}}\leq D\|w\|_{H^{s+2}}$. Choose
\begin{equation}\label{eq:cuspidal-contraction-radius}
 0<R\leq\min\{1,(4AC)^{-1},(4D)^{-1}\},\qquad \delta=R/(2A).
\end{equation}
The map $w\mapsto L_\varepsilon^{-1}(g-Q(w))$ takes the closed radius-$R$ ball in $H^{s+2}_0$ into its radius-$3R/4$ ball and has Lipschitz constant at most $1/2$. Its fixed point proves existence and uniqueness in that ball. Since $k\geq K_0/2$ and $\|\ddc w\|_{\mathrm{op}}\leq1/4$, the form is at least $K_0/4$. The embedding $H^{s+2}\subset C^2$ follows from $s\geq n+3$.
\end{proof}

The data vary in all torus variables. For example, writing $z_j=x_j+\mathrm i y_j$, a sufficiently small nonzero multiple of $\prod_j\cos(2\pi x_j)+\prod_j\cos(2\pi y_j)$ is an admissible zero-mean perturbation. The construction does not impose a separated form on the solution.

\begin{lemma}\label{lem:cuspidal-bounded-mass}
Every bounded $U\in\PSH(Y,k)$ has $\int_Y(k+\ddc U)^n=M$. The same holds with $H$ in place of $k$.
\end{lemma}
\begin{proof}
The products are those of \cref{prop:bounded-products}. The bounded local recursion and compact distributional Stokes give, for $r\geq1$,
\[
 \int_Y k^{n-r}\wedge(\ddc U)^r
     =\int_YU\,\ddc(k^{n-r})\wedge(\ddc U)^{r-1}=0.
\]
One may also convolve on the torus, add a fixed local positive quadratic, and pass through the bounded decreasing expansion; every smooth integral vanishes by the displayed identity. Expanding the top power proves the assertion. The proof for $H$ is identical, since its powers are also $\ddc$-closed and its mass is $M$.
\end{proof}

\begin{proposition}\label{prop:cuspidal-bounded-uniqueness}\index{Uniqueness!bounded nonclosed cuspidal equations}
Let $V$ be the solution in \cref{thm:cuspidal-quantitative-equation}. If any bounded $U\in\PSH(Y,k)$ satisfies $(k+\ddc U)^n=(k+\ddc V)^n$, then $U-V$ is constant.
\end{proposition}
\begin{proof}
Convolve all three objects with the same nonnegative radial torus kernel of radius $r_j\downarrow0$. Write
\[
 U_j=U*\varrho_j,\quad V_j=V*\varrho_j,\quad k_j=k*\varrho_j,
 \quad A_j=k_j+\ddc U_j,\quad B_j=k_j+\ddc V_j.
\]
Translation preserves positivity, so $A_j\geq0$ and $B_j\geq\beta K_0$, $\beta=1/4$. The forms $k_j$ converge smoothly and $V_j$ converges in $C^2$. On lifted coordinate balls choose a common positive quadratic $q$ with $\ddc q\geq k$. Convolution gives $q*\varrho_j=q+b r_j^2$ for a fixed $b>0$. The radial averages of the bounded psh functions $U+q$ and $V+q$ decrease as the radii decrease. Thus $U_j+b r_j^2+q$ and $V_j+b r_j^2+q$ decrease to the canonical values of $U+q$ and $V+q$. The identical constants cancel in their difference.

Put $\phi_j=U_j-V_j$ and $T_j=\sum_{r=0}^{n-1}A_j^r\wedge B_j^{n-1-r}$. The identities $\ddc k_j=0$ and $\mathrm dk_j\wedge\dc k_j=0$ persist under convolution: its torsion is still pulled back from $E_1\times E_2$. Differentiating the monomials therefore gives $\ddc T_j=0$. Also $T_j\geq B_j^{n-1}\geq\beta^{n-1}K_0^{n-1}$. Stokes applied to $\ddc(\phi_j^2)$ yields
\begin{equation}\label{eq:cuspidal-rough-energy}
 -\int_Y\phi_j(A_j^n-B_j^n)
 =\int_Y\mathrm d\phi_j\wedge\dc\phi_j\wedge T_j
 \geq\beta^{n-1}\int_Y\mathrm d\phi_j\wedge\dc\phi_j\wedge K_0^{n-1}.
\end{equation}
Only $\ddc$-closedness of $T_j$ is needed here.

We justify the four terms on the left separately. Expand $A_j^n$ relative to the local psh potential $U_j+b r_j^2+q$ and the smoothly converging coefficient $k_j-\ddc q$. The weighted limit \cref{lem:weighted-decreasing-products} gives $\int U_jA_j^n\to\int U(k+\ddc U)^n$. Ordinary decreasing convergence and uniform convergence of $V_j$ give the corresponding limit for $V_jA_j^n$. Since $B_j^n$ has uniformly converging continuous densities, bounded pointwise convergence and dominated convergence give the limit of $U_jB_j^n$. The term $V_jB_j^n$ converges by $C^2$ convergence. A finite partition of unity supplies the global statements. Equality of the two limiting measures makes the left side of \eqref{eq:cuspidal-rough-energy} tend to zero.

The flat Fourier Poincar\'e inequality now makes the $L^2$ norm of $\phi_j$ minus its average tend to zero. Bounded convergence gives $\phi_j\to U-V$ in $L^2$, so $U-V$ is constant almost everywhere. Local smooth corrections and the canonical psh representative extend this equality to every point.
\end{proof}

\begin{theorem}\label{thm:nonclosed-cuspidal-downstairs}
Let $f\in H^s(Y)$ be strictly positive and set
\[
 \nu=\pi_*(f\Omega),\qquad c=\frac{M}{\int_Y f\Omega}.
\]
If $\|cf-f_\varepsilon\|_{H^s}\leq\delta$, there is a unique bounded $u\in\PSH(X,h)$ with $\sup_Xu=0$ and
\[
 \mathrm{MA}^{\mathrm{ext}}_{X,h}(u)=c\nu.
\]
The function is continuous and locally H\"older of exponent $1/3$ in ambient coordinates. The measure $\nu$ is non-pluripolar and gives the singular locus zero mass. On the regular locus its density relative to $h^n$ is $f/J$, where $J=H^n/\Omega$; near the cusp it is $\lambda f/\rho$. The equation holds on every resolution and every common smooth refinement. The scalar is forced by the datum for every bounded solution.
\end{theorem}
\begin{proof}
The function $g=cf-f_\varepsilon$ has zero mean. Take the solution $w$ from \cref{thm:cuspidal-quantitative-equation}. The potential $v_0+w$ is bounded and $H$-psh. Descent through the resolution $\pi$ in \cref{thm:hermitian-space-descent} gives
\[
 u=(v_0+w)\circ\pi^{-1}-\sup_Y(v_0+w).
\]
This retains the exact disk background. The normalization is a homeomorphism, so $u$ is continuous. The extended-measure definition and \cref{thm:space-extended-measures} give the equation. Any bounded competing solution pulls back to a bounded $H$-psh potential. Since $\pi$ is a Borel homeomorphism, equality of the pushforward measures is equality upstairs. Subtract $v_0$ and apply \cref{prop:cuspidal-bounded-uniqueness}. The mass identity \cref{lem:cuspidal-bounded-mass} forces $c=M/\nu(X)$ before normalization.

The singular locus pulls back to the smooth-volume-null divisor $\{p\}\times E_2\times\cdots\times E_n$. A local pluripolar witness pulls back to a nontrivial psh witness away from that divisor, whose pole set has smooth volume zero. Thus $\nu$ is non-pluripolar. On the regular locus change of variables gives the density $f/J$. Near the cusp $\alpha=0$, so
\[
 H^n=n!\rho\,\kappa_1\wedge\cdots\wedge\kappa_n,
 \qquad \Omega=n!\lambda\,\kappa_1\wedge\cdots\wedge\kappa_n.
\]
The ratio is $\lambda f/\rho$; the degree and both factorials are accounted for.

For the H\"older assertion put $d=z-z'$ and $a=z+z'$. Then
\[
 z^2-z'^2=da,\qquad z^3-z'^3=d(3a^2+d^2)/4.
\]
If $|a|\geq|d|/4$, the first modulus is at least $|d|^2/4$, which controls $|d|^3$ on a bounded chart. Otherwise the second is at least $13|d|^3/64$. Hence
$|z-z'|\leq C(|z^2-z'^2|+|z^3-z'^3|)^{1/3}$. The $C^1$ regularity of $v_0+w$ and the other product coordinates give the asserted ambient estimate.

Every resolution $\sigma:Z\to X$ factors through $Y$: in the cusp chart the quotient $y/x$ is bounded off its exceptional set and extends on the normal smooth source, giving the normalization coordinate. Write $\sigma=\pi\circ\mu$. The map $\mu$ is a proper modification. The bounded pullback products agree off its exceptional locus and give that analytic locus zero mass by \cref{prop:bounded-products}. Their pushforwards therefore agree on all of $Y$. This proves the equation on every resolution, and the same argument applies to further common refinements. The full potential is literally $u\circ\sigma$, with its finite constant value on every fiber; surjectivity preserves all ordinary comparisons.
\end{proof}

Existence alone has a larger range of data. The positive repaired reference puts the normalization in the bounded-mass regime of \cref{chap:bounded-mass-equations}. The smallness condition above supplies quantitative regularity and a positive solution needed for the uniqueness comparison.

\Needspace{24\baselineskip}
\begin{corollary}\label{cor:cuspidal-lp-bounded-existence}
For the same $X,h$, every nonzero nonnegative $f\in L^p(Y,\Omega)$, $p>1$, admits a bounded normalized $u\in\PSH(X,h)$ with
\[
 \mathrm{MA}^{\mathrm{ext}}_{X,h}(u)=\frac{M}{\int_Yf\Omega}\,\pi_*(f\Omega).
\]
\end{corollary}
\begin{proof}
The torus $Y$ is K\"ahler and has bounded mass by \cref{lem:kahler-modification-bounded-mass}. Since $k>0$, the zero function is its ordinary model reference: $P_k[0]=0$. Its lower reference volume is $M>0$ by \cref{lem:cuspidal-bounded-mass}, because the ordinary type of zero consists precisely of bounded potentials. Thus \cref{thm:li-general-measure} gives a bounded $k$-psh solution for $f\Omega$, and its scalar is $M/\int f\Omega$. Add $v_0$ and use \cref{thm:hermitian-space-descent,thm:space-extended-measures} to obtain the stated normalized equation downstairs. Non-pluripolarity of the datum follows as in \cref{thm:nonclosed-cuspidal-downstairs}. No smallness or strict positivity of $f$ is needed for this existence assertion.
\end{proof}

These spaces have no rational curves. A map $\mathbb P^1\to X$ lifts locally to $Y$ by the cusp normalization, and the lifts glue by uniqueness. Every pulled-back flat holomorphic one-form vanishes on $\mathbb P^1$, so the lifted map has zero derivative and is constant. The results retain this explicit geometry, bounded ordinary type and power-closed torsion. Ambient H\"older regularity and uniqueness throughout the bounded class are proved for the small normalized Sobolev neighborhood. The general equation problem \cref{qu:singular-space-equations-program}, including arbitrary pole types and nonclosed torsion, remains open.
